\documentclass[10pt,letterpaper]{article}

% ---------- ATS-SAFE PACKAGES ONLY ----------
\usepackage[letterpaper,margin=0.55in,top=0.25in,bottom=0.25in]{geometry}
\usepackage[T1]{fontenc}
\usepackage[utf8]{inputenc}
\usepackage{enumitem}
\usepackage{titlesec}
\usepackage{hyperref}
\usepackage{xcolor}
\usepackage{parskip}
\usepackage{helvet}
\renewcommand{\familydefault}{\sfdefault}

% Standard, universally-parsed sans font (no icon fonts, no color boxes, no tables)
\definecolor{linkblue}{RGB}{0,71,171}
\hypersetup{
    colorlinks=true,
    linkcolor=linkblue,
    urlcolor=linkblue,
    pdfborder={0 0 0}
}

\pagestyle{empty}
\setlength{\parindent}{0pt}

% Section formatting: plain bold caps + a single rule. No colorbox, no shading, no tables.
\titleformat{\section}{\bfseries\large}{}{0em}{}[\vspace{2pt}\hrule]
\titlespacing{\section}{0pt}{3pt}{2pt}
\setlength{\parskip}{1pt}

\newcommand{\resitem}[1]{\item #1}
\setlist[itemize]{topsep=0pt, itemsep=0pt, parsep=0pt, leftmargin=16pt, label=\textbullet}

% Simple two-column-look entry using a plain \hfill (still linear text order for ATS)
\newcommand{\entry}[4]{%
    \textbf{#1} \hfill \textbf{#3}\\
    \textit{#2} \hfill \textit{#4}\\[0pt]
}

\begin{document}

% ---------- HEADER (plain text, centered, no image/logo) ----------
\begin{center}
    {\LARGE \textbf{Abhishek Raj}}\\[4pt]
    New Delhi, India \quad $\vert$ \quad +91-8810590060 \quad $\vert$ \quad \href{mailto:abhishekraj185a@gmail.com}{abhishekraj185a@gmail.com}\\[2pt]
    \href{https://www.linkedin.com/in/abhishekraj0011/}{LinkedIn} \quad $\vert$ \quad
    \href{https://github.com/i-abhishek-04}{GitHub} \quad $\vert$ \quad
    \href{https://portfolio-theta-blue-ykgevsy0wm.vercel.app/}{Portfolio} \quad $\vert$ \quad
    \href{https://leetcode.com/u/iAbHi__004/}{LeetCode}
\end{center}

\vspace{-4pt}

% ---------- SUMMARY (optional, boosts keyword density for ATS) ----------
\section*{Summary}
Geoinformatics undergraduate skilled in full-stack development (Python, JavaScript, React, FastAPI) and applied AI/LLM engineering (LangChain, RAG, prompt engineering), with hands-on experience shipping tested, API-integrated projects end to end.

% ---------- EXPERIENCE ----------
\section*{Experience}

\entry{Frontend Web Development Intern}{Magnus Corps (Remote)}{May 2026 -- Present}{}
\vspace{-4pt}
\begin{itemize}
    \item Built responsive frontend layouts with semantic HTML/CSS and JavaScript-driven interactivity for the agency's client-facing marketing site.
    \item Helped integrate an AI-powered feature into the frontend and maintained a clear, reusable component structure aligned with the design system.
    \item Tech stack: HTML, CSS, JavaScript, AI integration.
\end{itemize}

% ---------- PROJECTS ----------
\section*{Projects}

\entry{SyncPad --- Real-Time Collaborative Code Editor (Hand-Built CRDT)}{\href{https://github.com/i-abhishek-04/SyncPad}{GitHub} $\vert$ \href{https://sync-pad-rho.vercel.app}{Live Demo}}{2026}{}
\vspace{-4pt}
\begin{itemize}
    \item Engineered a real-time collaborative editor from scratch with a hand-built RGA CRDT engine (Roh et al. 2011), achieving strong eventual consistency across concurrent clients without server-side locking.
    \item Built a FastAPI and native WebSocket backend with a React/Vite frontend featuring live cursor presence and a CRDT inspector; validated with an 8-test convergence suite covering concurrent edits, 4-client convergence, and reconnect resync.
    \item Tech stack: Python, FastAPI, WebSockets, React, Vite, PrismJS, CRDT/RGA.
\end{itemize}

\entry{OmniSocial AI --- Full-Stack Social Media Analytics SaaS}{\href{https://github.com/i-abhishek-04/OmniSocial-AI.git}{GitHub} $\vert$ \href{https://omnisocial-ai.vercel.app/}{Live Demo}}{2026}{}
\vspace{-4pt}
\begin{itemize}
    \item Architected a full-stack creator analytics platform with FastAPI and React 19, exposing 21 REST endpoints backed by JWT authentication and 4 SQLAlchemy models, validated by a 33-test pytest suite (100\% passing); integrated a pluggable platform-adapter system across 5 live third-party APIs (YouTube, GitHub, Instagram, Reddit, Dev.to) with automatic fallback.
    \item Integrated an LLM-powered assistant with tiered fallback across Groq (Llama 3.3 70B) and Google Gemini, grounding responses in live user analytics; conducted a security audit hardening JWT fallback, rate limiting, and CORS handling.
    \item Tech stack: Python, FastAPI, React 19, Vite, SQLAlchemy, JWT, SQLite, pytest, Groq, Gemini.
\end{itemize}

\entry{LangChain Multi-Provider AI Studio \& Chain Playground}{\href{https://github.com/i-abhishek-04/Langchain-enterprise-rag-platform.git}{GitHub} $\vert$ \href{https://langchain-enterprise-ai-platform.vercel.app}{Live Demo}}{2026}{}
\vspace{-4pt}
\begin{itemize}
    \item Built an interactive React/Vite studio unifying 5 free-tier AI providers (Groq, Gemini, OpenRouter, HuggingFace, local Ollama) behind one interface, with a zero-config client-side fallback mode and an LCEL Chain Flow Visualizer.
    \item Curated a prompt template library (ReAct, RAG Q\&A, JSON extraction) and a Python code exporter for one-click runnable LangChain snippets, structured within a uv-managed Python monorepo.
    \item Tech stack: React, Vite, Python, LangChain, LCEL, Groq, Gemini API, HuggingFace, Ollama.
\end{itemize}

% ---------- TECHNICAL SKILLS ----------
\section*{Technical Skills}
\textbf{Programming Languages:} C++, Python, JavaScript, HTML, CSS, SQL\\
\textbf{Frameworks \& Tools:} React, Vite, FastAPI, SQLAlchemy, Git/GitHub, REST APIs, MySQL, MongoDB\\
\textbf{Generative AI / LLM Tools:} LangChain, LCEL, Google Gemini API, Groq, Ollama, OpenRouter, Prompt Engineering, RAG (Retrieval-Augmented Generation)\\
\textbf{Core CS Subjects:} Data Structures \& Algorithms, Object-Oriented Programming, Operating Systems, Computer Networks, Database Management Systems, System Design

% ---------- EDUCATION ----------
\section*{Education}

\entry{Netaji Subhas University of Technology}{B.Tech, Geoinformatics}{2023 -- 2027}{}
CGPA: 6.92 / 10

\vspace{3pt}
{\small Class XII (CBSE), 2023 \quad $\vert$ \quad Class X (CBSE), 2021}

% ---------- ACHIEVEMENTS ----------
\section*{Achievements}
\begin{itemize}
    \item Solved 400+ Data Structures \& Algorithms problems on LeetCode spanning graphs, dynamic programming, and trees.
    \item Ranked within the top 20 percentile in JEE Main among 1M+ candidates.
\end{itemize}

\end{document}